% Reference snapshot from revised nthm.tex, 18 October 2026.
% This excerpt uses the author's manuscript macros and is not an extra input.
% It is stored here so the experimental hypotheses can be read on the branch
% without recovering conversation attachments. It introduces no new Lean axiom.
\begin{definition}
	Let $\Stree$ be an $\S$-tree, $F\colon T\to T$ a shape-preserving function, and $m\geq0$ an integer. We say that $F$ \emph{skips level $m$} if $m\notin\ell[F[T]]$, and that $F$ \emph{skips only level $m$} if $\ell[F[T]]=\omega\setminus\{m\}$.
\end{definition}
A shape-preserving function skipping only level $m$ fixes $T({<}m)$.
It sends level $m$ to level $m+1$ and has consecutive image levels
thereafter. Thus it is determined by its restriction to $T(m)$.

\begin{definition}[Successor monoid]\label{def:successor-monoid}
Given an $\S$-tree $\Stree$ and a set $\M$ of shape-preserving functions
$T\to T$, we call $\SNtree$ a \emph{successor monoid} if the following
condition holds.
\begin{enumerate}[label=(M\arabic*)]
		\item\label{item:Nmonoid}
		\emph{The set $\M$ forms a closed monoid}: $\M$ contains the identity function $\mathrm{Id}\colon T\to T$,
		for every $F,F'\in \M$ it holds that $F\circ F'\in \M$, and every sequence $(F_i)_{i\in \omega}$ of functions in $\M$
		satisfying
		$(\forall_{i\in \omega})F_i\restriction_{T({\leq}i)}=F_{i+1}\restriction_{T({\leq}i)}$ has a limit in $\M$, that is, a function $F_\infty \in \M$ such that $(\forall_{i\in \omega})F_\infty\restriction_{T({\leq}i)} = F_i\restriction_{T({\leq}i)}$. Equivalently, one can define the \emph{pointwise convergence topology} on the set of shape-preserving functions in the typical way, and then we demand that $\M$ is closed in this topology.
\end{enumerate}
\end{definition}
The monoid and its pointwise topology will be kept throughout. Blocks
are restrictions of its functions, not replacements for the functions
or extra annotations on them. We use the following notation:

$$\begin{aligned}
 \M^n&=\{F\in\M:F\restriction_{T({<}n)}\text{ is the identity}\},\\
 \AM&=\{F\restriction_{T({<}k)}:F\in\M,\ k\in\omega\},\\
 \AM^n_k&=\{F\restriction_{T({<}n+k)}:F\in\M^n\}.
\end{aligned}$$
Thus $\AM^n_k$ consists of realised approximations with $n$ levels
fixed and $k$ levels moving. For $n+k>0$, write $\AM^n_k(m)$ for
those whose last image level is $m$. The empty approximation belongs
to $\AM$.

We separate the local operations needed for completed blocks from the
stronger one-level operations used in the applications. To describe a
line over a finite alphabet $\Gamma$, take a word $L$ over
$\Gamma\cup\{\lambda\}$ in which $\lambda$ occurs. Write $L(e)$ for
substitution of $e$ at every occurrence of $\lambda$, and $L(*)$ for
the prefix before its first occurrence. We write $u\conc v$ for
concatenation.

\begin{definition}[$(\S,\M)$-tree]\label{def:block-monoid}
Let $\SNtree$ be a successor monoid with every level inhabited.
Fix finite alphabets $\Gamma_p\subseteq\AM^p_1(p+1)$ for $p\in\omega$.
We call it an \emph{$(\S,\M)$-tree}, relative to these alphabets,
if the following conditions hold.
\begin{enumerate}[label=(M\arabic*)]
\setcounter{enumi}{1}
\item\label{item:Mcompletion}\emph{Canonical completion.}
Every $f\in\AM$ with domain $T({\leq}n)$ has an extension $f^+\in\M$ such that
$$
 \widetilde{f^+}(n+i)=\widetilde f(n)+i\qquad(i\geq0).
$$
We put $\emptyset^+=\mathrm{Id}$.

\item\label{item:Mendpoint}\emph{Block factorisation.}
If $f,g\in\AM$ and $g$ extends $f$ by one source level, then
$$
 g^+=D\circ f^+\qquad\text{for some }D\in\M.
$$

\item\label{item:Mseeded}\emph{Realisation after a recording prefix.}
For $w=e_0\cdots e_{k-1}\in\Gamma_p^{<\omega}$ put
$$
 g_w=(e_0^+\circ\cdots\circ e_{k-1}^+)
           \restriction_{T({\leq}p)},
$$
where the empty product is $\mathrm{Id}$.
There is $s_p\in\Gamma_p^{<\omega}$ such that every line $L$ over
$\Gamma_p$ admits one $h\in\AM^p_2$ with
$$
 h\restriction_{T({\leq}p)}=g_{s_p\conc L(*)},\qquad
 h\circ e=g_{s_p\conc L(e)}\quad(e\in\Gamma_p).
$$
\end{enumerate}
\end{definition}
Canonical completions are unique by the successor equation; see
Proposition~\ref{prop:canonical}. In~\ref{item:Mendpoint}, $D$ acts
\emph{after} $f^+$; this is not the right-composition reduction used
below. The fixed prefix and exact endpoints of $D$ follow automatically
(Proposition~\ref{prop:block-extraction}). Iterating gives the same
factorisation for every finite extension $f\subseteq g$.

The recording word is fixed before the line is chosen, and a single
$h$ works for every letter. When $\Gamma_p=\emptyset$, take $s_p=()$
and $h$ the identity restriction, so no additional condition is imposed.
The alphabets are designated families, not necessarily all one-level
maps. For ordered words they consist of constant appendings; the strong
axioms below permit the full alphabet. Proposition~\ref{prop:block-local-equivalence}
shows that the completion--factorisation pair is equivalent to the
local completion--extraction formulation: no extra uniform choice of
factors or completion of arbitrary partial maps is required.

\begin{definition}[Strong $(\S,\M)$-tree]\label{def:sntree}
A successor monoid $\SNtree$ is a \emph{strong $(\S,\M)$-tree} if it
satisfies the following two additional conditions. These are the
original decomposition and duplication axioms, now distinguished by
primes.
\begin{enumerate}[label={(M\arabic*$'$)}]
\setcounter{enumi}{1}
		\item\label{item:Ndecomposition} \emph{The set $\M$ admits decompositions}:
			For every $n\in \omega$ and $F\in \M$ such that $\tilde{F}(n)>0$ and $F$ skips level $\tilde{F}(n)-1$,
		there exist $F_1, F_2\in \M$ such that $F_2$ skips only level $\tilde{F}(n)-1$ and $$F_2\circ F_1\restriction_{T({\leq} n)}=F\restriction_{T({\leq} n)}.$$

		\item\label{item:Nduplication} \emph{The set $\M$ is closed under duplication}: For all $n$ and $m$ with $n<m\in \omega$, there exists a function $F_m^n\in \M$
		skipping only level $m$ such that
		for every $a\in T(n)$, $b\in T(m)$, $\bar{p}\in T^{<\omega}$, and $c\in \Sigma$, where $\S(a,\bar{p},c)$ is defined and $\S(a,\bar{p},c)\preceq b$,
		we have $F^n_m(b)=\S(b,\bar{p},c)$.
\end{enumerate}
\end{definition}
